\section{讨论：证据强度与开放问题}

直接证据最充分的是 BaZn$_2$Si$_2$O$_7$ 的相变、Ba/Sr 固溶和 Zn 位替位；高温 XRD、单晶 XRD 与膨胀测量在这些研究中相互印证 \cite{lin1999phase,thieme2015ba1,kerstan2012thermal,thieme2016new}。玻璃析晶方面，成核剂和微结构效应已有系统实验，但不同研究的玻璃组成、升温速率和晶粒尺寸并不统一，导致“相稳定性”与“样品宏观膨胀”不能混为一谈 \cite{thieme2016thermal,vladislavova2018crystallisation}。

第一项关键缺口是 BaO--ZnO--SiO$_2$ 三元相图和热力学函数。当前只能借助二元相平衡、实验室 XRD 和热分析拼接相区；理论工作提示必须纳入振动自由能和温度效应 \cite{erlebach2021phase}。第二项缺口是 Ba$_5$Y$_{12}$Zn[O(SiO$_4$)]$_8$ 的直接结构证据。仅凭“钡—钇—锌—硅酸根”共存不能建立同构关系，需获得可复核的衍射数据和配位精修。

建议的安全、可证伪实验路径是：先在小规模、封闭式高温炉和经审查的坩埚材料条件下制备组成梯度；同步采集室温/高温粉末 XRD、TG-DSC、Raman 和显微成分图；对出现单一衍射相的样品再做同步辐射或单晶结构解析。该建议沿用 BaZn$_2$Si$_2$O$_7$ 研究中 HT-XRD 与膨胀测量的互补思路 \cite{kerstan2012thermal,thieme2016very}，不把未验证的 Y-rich 配方当作可直接执行的实验协议。

结构描述符可帮助缩小候选范围。负热膨胀研究表明，四面体转动、链柔性和配位多面体空隙可作为筛选特征，但描述符只能提出候选，不能替代相平衡实验 \cite{kireeva2022pauling,erlebach2021phase}。
